% ============================================================
%  第一章  引言
% ============================================================
\section{引言}

\subsection{编写目的}

本文档是“鹰眼长空”（低空无人机智能监测系统，版本 V1.0）的程序设计说明书，用于完整描述该软件的系统架构、程序流程、模块划分、核心算法、数据结构与接口约定，为软件的开发维护、功能扩展、测试验收以及软件著作权登记提供技术依据。

文档的预期读者包括：项目开发人员、测试人员、系统集成人员、技术评审人员以及后续承接本项目的维护人员。阅读本文档要求读者具备 Python 编程、Web 前后端开发、计算机视觉（目标检测与双目立体视觉）和三维数据可视化的基础知识。

具体而言，本文档承担以下职责：
\begin{enumerate}
  \item 说明软件的建设目标、功能边界与运行环境，明确“做什么”与“不做什么”；
  \item 给出系统的分层架构、模块划分、模块间依赖关系与进程部署视图；
  \item 用流程图描述系统启动、数据处理、检测推理、轨迹合成、备份恢复等关键流程；
  \item 逐一说明各功能模块的类、函数、算法公式与异常处理策略；
  \item 定义轨迹数据文件、内存数据结构、配置项和 REST 接口的格式；
  \item 说明当前版本中尚未完成的轨迹预测功能的技术方案与开发计划。
\end{enumerate}

\begin{notebox}[阅读约定]
\begin{itemize}
  \item 文中路径均相对于项目根目录 \file{drone/}；代码标识符以等宽字体表示。
  \item 流程图符号：矩形为处理，菱形为判断，圆角矩形为起止，平行四边形为输入输出；实线为控制流，虚线为数据流或异步通知。
  \item 除特别说明外，坐标单位为米（m），时间单位为秒（s），角度单位为度（$^{\circ}$）。
  \item 第 10 章截图取自 V1.0 实际运行界面（默认深色主题，1600$\times$1000 视口），图中圆形数字标记与图题、正文一一对应。
\end{itemize}
\end{notebox}

\subsection{术语与缩略语}

本文中出现的专业术语与缩略语如表~\ref{tab:terms} 所示。

\begin{table}[H]
\centering
\caption{术语与缩略语}
\label{tab:terms}
\small
\begin{tabularx}{\linewidth}{@{}>{\heiti}p{3.4cm}X@{}}
\toprule
术语 / 缩写 & 含义说明 \\
\midrule
YOLO & You Only Look Once，一类单阶段实时目标检测神经网络模型 \\
ultralytics & YOLOv8/YOLO11 系列模型的官方 Python 实现与训练框架 \\
双目立体视觉 & 利用两个相机对同一目标的视差恢复目标深度的技术 \\
基线（Baseline） & 双目相机两个光心之间的物理距离，记作 $B$，单位米 \\
视差（Disparity） & 同一目标在左右目图像上的横坐标之差，记作 $d$，单位像素 \\
三角测量 & 由视差与焦距、基线按几何关系解算目标三维坐标的过程 \\
ByteTrack / BoT-SORT & 两种多目标跟踪算法，由 ultralytics 内置支持 \\
轨迹重建 & 将各传感平台采集的离散轨迹点整合为统一时空坐标系的连续轨迹 \\
综合轨迹 & 按权重对多平台轨迹做加权平均得到的融合轨迹（synthetic） \\
线性外推 & 以最后两点的速度向量延长轨迹的预测方法，本系统作为兜底方案 \\
捕捉时机 & 目标高度低、速度小、加速度平稳、曲率小的时刻，便于实施反制 \\
REST API & 基于 HTTP 与 JSON 的表述性状态转移接口 \\
蓝图（Blueprint） & Flask 中用于组织一组路由的模块化机制 \\
pywebview & 使用系统原生 WebView 承载网页的轻量桌面窗口库 \\
Three.js & 基于 WebGL 的浏览器三维渲染库 \\
ECharts & 百度开源的数据可视化图表库 \\
\bottomrule
\end{tabularx}
\end{table}

\subsection{参考资料}

\begin{enumerate}
  \item 《鹰眼长空 — 用户使用手册》（\file{docs/manual.md}）
  \item 《鹰眼长空 — 配置说明》（\file{docs/CONFIG.md}）
  \item 《鹰眼长空 — 开发文档》（\file{docs/DEVELOPER.md}）
  \item 《Phy-ODE-Diffusion 轨迹预测模型说明》（\file{docs/hybrid_prediction_model.md}）
  \item GB/T 8567—2006《计算机软件文档编制规范》
  \item GB/T 28169—2011《嵌入式软件 通用要求》（参考其文档组织方式）
  \item ultralytics 官方文档（\url{https://docs.ultralytics.com}）
  \item Three.js 官方文档（\url{https://threejs.org/docs}）
\end{enumerate}

\subsection{文档结构}

全文共十二章，另附三个附录，各章内容安排如下：

\begin{table}[H]
\centering
\caption{文档结构说明}
\label{tab:docstructure}
\small
\begin{tabularx}{\linewidth}{@{}>{\heiti}p{2.6cm}X@{}}
\toprule
章节 & 主要内容 \\
\midrule
第 1 章 & 编写目的、术语、参考资料与文档结构 \\
第 2 章 & 软件概述：建设目标、功能清单、运行环境、技术栈与软件组成 \\
第 3 章 & 系统总体设计：设计原则、分层架构、模块划分、依赖与部署视图 \\
第 4 章 & 程序流程设计：启动、请求、检测、合成、预测、备份等流程与状态机 \\
第 5 章 & 轨迹重建与分析模块的详细设计 \\
第 6 章 & 无人机检测模块的详细设计 \\
第 7 章 & 轨迹预测功能的现状说明与后续开发方案 \\
第 8 章 & 数据结构与文件格式设计 \\
第 9 章 & 接口设计：命令行、页面路由、REST 接口、外部 AI 接口 \\
第 10 章 & 软件使用说明：界面风格、逐页操作、标注截图、典型流程与常见问题 \\
第 11 章 & 安装、部署与运行说明 \\
第 12 章 & 测试与验证方案 \\
附录 & 源文件清单、配置项全表、版本记录 \\
\bottomrule
\end{tabularx}
\end{table}
